\documentclass[11pt]{article}
\usepackage[a4paper,margin=23mm]{geometry}
\usepackage[T1]{fontenc}
\usepackage{lmodern,amsmath,amssymb,amsthm,booktabs,longtable,array,microtype}
\usepackage{xcolor}
\usepackage[colorlinks=true,linkcolor=blue!50!black,urlcolor=blue!50!black,citecolor=blue!50!black]{hyperref}
\usepackage{enumitem}
\setlist{nosep,leftmargin=*}
\newtheorem{theorem}{Theorem}
\newtheorem{lemma}[theorem]{Lemma}
\newcommand{\R}[1]{\texttt{R12\_1\_#1}}
\newcommand{\FPdim}{\operatorname{FPdim}}
\renewcommand{\Vec}{\operatorname{Vec}}
\newcommand{\Rep}{\operatorname{Rep}}
\newcommand{\Irr}{\operatorname{Irr}}
\newcommand{\C}{\mathbb C}
\newcommand{\Z}{\mathbb Z}
\newcommand{\Q}{\mathbb Q}
\setlength{\parskip}{4pt}
\setlength{\parindent}{0pt}
\title{Rank twelve, multiplicity one:\\third offline theoretical continuation}
\author{Fusion-ring census research notes}
\date{26 October 2026}
\begin{document}
\maketitle
\begin{abstract}
We revisit all 66 rings left unresolved by the previous round. A unified
Morita and cohomology argument resolves the fifteen dimension-120 cases,
giving three unitary models and twelve general-complex exclusions.
Twisted Ising constructions give two further models; relative cocycle
adaptation gives another two. A complete central-induction and indicator
certificate excludes one quadratic ring. Thus this round resolves twenty
cases: thirteen exclusions and seven models. Of the 632 supplied rings,
438 are excluded, 148 have unitary models, and 46 remain unresolved.
All results are offline. Independent computational and mathematical checks
within the project have been performed; independent human verification
remains pending.
\end{abstract}

\section{Results and precise scope}
The objects counted here are supplied based fusion rings, identified by
their permanent census IDs and SHA256 hashes of their complete tensors.
These are not counts of fusion categories up to tensor equivalence.
An exclusion concerns every complex fusion categorification. A positive
result is an actual unitary category with all $12^3=1728$ structure
constants matched under a unit-preserving basis bijection.

\begin{center}
\begin{tabular}{lrrr}
\toprule
Stage & Excluded & Modeled & Unresolved\\
\midrule
Previous completed round &425&141&66\\
This round: changes &+13&+7&$-20$\\
Current total &\textbf{438}&\textbf{148}&\textbf{46}\\
\bottomrule
\end{tabular}
\end{center}

The new exclusions are
\[
\begin{gathered}
\R{161},\ \R{211},\ \R{262},\ \R{265},\ \R{312},\ \R{350},\ \R{362},\\
\R{388},\ \R{453},\ \R{469},\ \R{472},\ \R{503},\ \R{516}.
\end{gathered}
\]
The new unitary models are
\[
\R{142},\ \R{240},\ \R{361},\ \R{408},\ \R{414},\ \R{444},\ \R{531}.
\]
No earlier exclusion or model is retracted. The previous round is preserved
as a frozen snapshot in the archive. Failed searches, feasible necessary
conditions, and based-ring automorphisms without categorical lifts have
not been used to change any status.

\begin{center}
\begin{tabular}{lrrrr}
\toprule
Family & Entering & Excluded here & Modeled here & Remaining\\
\midrule
Dimension 32 &20&0&2&18\\
Dimension 48 &22&0&2&20\\
Dimension 120 &15&12&3&0\\
Dimension 144 &4&0&0&4\\
Quadratic, $\sqrt{13}$ &1&0&0&1\\
Quadratic, $\sqrt{10}$ &4&1&0&3\\
\midrule
Total &66&13&7&46\\
\bottomrule
\end{tabular}
\end{center}

\section{The dimension-120 family is completely resolved}
\begin{theorem}
Among the sixteen supplied rings of type $(1^4,2^4,5^4)$, the categorifiable
ones are precisely $\R{25}$, $\R{142}$, $\R{361}$ and $\R{444}$.
Each has a unitary group-theoretical realization. The first was already
modeled; the other three and the twelve exclusions are new decisions.
\end{theorem}

\subsection{Two forced pointed Morita reductions}
Every source has a dimension-20 subring of type $(1^4,2^4)$. There is a
unique common involution $b$ fixing the four dimension-two objects on both
sides. On the dimension-five objects it acts freely, with identical left
and right permutations. These are exact properties of all sixteen tensors.

A fixed simple trivializes the pointed associator on $\langle b\rangle$,
so $A=1+b$ is an actual separable algebra. The cyclic stabilizer argument
from the previous round makes the dimension-20 subcategory pointed under
$A$-bimodule duality. Briefly, a fixed dimension-two simple has an $A$-module
structure because $H^2(C_2,\C^\times)=0$. Its internal end has dimension two
and contains the unit; hence it is supported on $1+b$. The same vanishing
supplies a commuting left action. The free dimension-five orbits each
give two dimension-five bimodules. Thus the full dual $\mathcal D$ has type
$(1^{20},5^4)$.

The invertible group of $\mathcal D$ has a unique Sylow subgroup
$H=C_5$. It acts trivially on the four other simples on both sides.
The algebra $B=\bigoplus_{h\in H}h$ therefore exists. For a simple right
$B$-module with underlying dimension-five object, the internal end has
dimension $25/5=5$ and contains the unit. It cannot contain a noninvertible
simple. Its invertible support is a subgroup of order five and hence
equals $H$. This proves that a commuting left action exists; it does not
assume that the original object's full stabilizer had order five.

Each of the four dimension-five objects gives $5\cdot5$ invertible
bimodules. The pointed sector gives twenty more. Consequently
$\mathcal E={}_B\mathcal D_B$ is pointed of dimension 120. This proves
group-theoreticity for any categorification, with no unitary or spherical
hypothesis.

The reverse algebra for the first dual is $R=1+c$, where $c$ has order two
and lies outside $H$. The actual free $B$--$R$ bimodule $B\otimes R$ has
ten distinct invertible constituents and transitive support, so is simple.
Its internal end in $\mathcal E$ has dimension
$10^2/(5\cdot2)=10$. The original category therefore has the form
\[
\mathcal C(G,\omega,K,\psi),\qquad |G|=120,\quad |K|=10.
\]
The subgroup order is obtained from an actual algebra, not guessed from
a dimension-divisibility condition.

\subsection{The group and subgroup are forced}
If $K=C_{10}$, the identity double coset supplies ten invertibles, which
is impossible. Therefore $K=D_{10}$, where the subscript denotes group
order. Its Schur multiplier is trivial. Each normalizer coset supplies
degrees $(1,1,2,2)$, so the normalizer has order twenty. The two remaining
double cosets have size fifty and stabilizer $C_2$; these supply the four
dimension-five objects.

If $G$ were solvable, a Hall subgroup of order forty containing $K$
would be a union of $K$-double cosets. It can contain none of size fifty,
and would have to lie in the order-twenty normalizer: a contradiction.
The nonsolvable groups of order 120 are $A_5\times C_2$, $S_5$ and
$\mathrm{SL}(2,5)$. The last has no $D_{10}$ subgroup. In $S_5$, or for
the trivial graph $D_{10}\subset A_5\times C_2$, an induced faithful
$C_2$-grading has neutral type $(1^2,2^2,5^2)$. A dimension-five simple
times its dual has dimension 25, larger than the sum $2+4+10=16$ of all
neutral simple dimensions. Such a product cannot be multiplicity-free.

Only the reflection-sign graph remains:
\[
G=A_5\times C_2,\qquad
K=\{(h,\operatorname{sgn}_{D}(h)):h\in D_{10}\}.
\]

\subsection{Exhaustive cohomology and the four outputs}
The admissible classes form the kernel of restriction from
$H^3(G,\C^\times)$ to $H^3(K,\C^\times)$.
Its 5-primary part vanishes because the subgroup index twelve is prime
to five. Its 3-primary part is $C_3$ and does not change the fusion ring,
by the lemma below.

The ambient 2-primary part is $C_2^3$. Restriction to
$V_4\times C_2$ is injective, and the normalizer $C_3$ has a
three-dimensional invariant subspace in its degree-three cohomology.
Three independent global classes attain this bound: $u$, the transfer
of the anisotropic $V_4$ cocycle; $v$, the external $C_2$ cube; and $w$,
the cup product of the binary-icosahedral Schur cocycle with the external
character. Each restricts to the same nonzero class on $K$. The relative
2-primary part is thus generated by $u+v$ and $w+v$.

Explicit cochain gauges make these two cocycles identically trivial on
$K$, so the subgroup algebra has cochain $\psi=1$. There is no further
subgroup-cochain parameter because $H^2(D_{10},\C^\times)=0$.

\begin{center}
\begin{tabular}{ccc}
\toprule
Coefficient of $u+v$ & Coefficient of $w+v$ & Exact tensor\\
\midrule
0&0&$\R{25}$\\
0&1&$\R{142}$\\
1&0&$\R{361}$\\
1&1&$\R{444}$\\
\bottomrule
\end{tabular}
\end{center}

\begin{lemma}[Coprime primary twists preserve the based fusion ring]
Let $p\nmid |K|$. A $p$-primary change of the associator of a
group-theoretical category $\mathcal C(G,\omega,K,\psi)$, with a compatible
$p$-primary subgroup trivialization, does not change its based fusion ring.
\end{lemma}
\begin{proof}
Choose finite roots of unity for the data and a complete strictly henselian
discrete valuation ring of mixed characteristic $(0,p)$ containing them.
The $p$-power roots reduce to one. Every double-coset stabilizer has order
dividing $|K|$, so its twisted group algebra is separable over this ring
and has split semisimple residue fiber. Its simple modules lift with the
same ranks. The balanced tensor product and the equivariant Hom spaces
are direct summands furnished by averaging idempotents whose denominators
divide $|K|$. Their ranks, hence all fusion coefficients, are unchanged
under reduction. Both associators have identical reductions.
The argument needs $p\nmid |K|$, not $p\nmid |G|$.
\end{proof}

An independent Mackey-formula proof is also included in the archive.
The four exact constructions, together with this categorical converse,
prove the twelve exclusions. A failed construction search alone would
not establish the theorem.

\section{Two twisted Ising models in dimension 32}
The same construction principle gives $\R{240}$ and $\R{408}$.
Let $\mathcal I$ be a unitary Ising category, with objects
$1,\psi,\sigma$ and $\sigma^2=1+\psi$. For a group $E$ of order sixteen,
a normalized unitary cocycle $\omega$, and a grading
$\epsilon:E\to C_2$, take the matched-grading fusion subcategory
\[
\mathcal B=\{(X,g)\in\mathcal I\boxtimes\Vec_E^\omega:
                  \deg X=\epsilon(g)\}.
\]
It has type $(1^{16},(\sqrt2)^8)$ and dimension 32.
If $s$ is an even involution and $\omega$ restricts trivially to
$\langle s\rangle$, then $A=1+(1,s)$ is a connected separable algebra.
The category ${}_A\mathcal B_A$ is an actual unitary Morita model.

For $\R{240}$, use $E=D_{16}$, with
\[
(a,b)(c,d)=([a+(-1)^bc]_8,b+d),\qquad \epsilon(a,b)=a\bmod2,
\]
and $s=(0,1)$. A suitable cocycle is
\[
\omega((a,b),(c,d),(e,f))=
\exp\!\left(\frac{2\pi i}{8}a(-1)^b
      \frac{c+(-1)^de-[c+(-1)^de]_8}{8}\right).
\]
It is the cup product of the crossed rotation coordinate with its twisted
integral Bockstein; the two sign coefficient systems cancel.
This proves the cocycle identity conceptually. All $16^4=65\,536$
identities are also checked exactly.

For $\R{408}$, replace the group by the semidihedral group
$E=\langle r,s:r^8=s^2=1,\ srs^{-1}=r^3\rangle$.
The supplied sixteenth-root cocycle is an exact lift of a normalized
modulo-eight cocycle. Its full table and all lift data are frozen; every
cocycle identity and the restriction to the chosen involution are
checked. This uses an actual pointed category and an actual subgroup
algebra, not a presumed lift of a fusion-ring symmetry.

Both bimodule categories have type $(1^8,2^2,(\sqrt8)^2)$.
The primary constructor computes exact balanced products; an independent
implementation computes full grade-preserving Hom spaces. Both match all
1728 coefficients to the original tensors. The Ising category supplies
the full associator, including its three-$\sigma$ matrix block, although
that block is not needed for these subgroup-bimodule operations.

\section{Two models from relative cocycle adaptation}
The following elementary device makes previously unusable cocycles usable.

\begin{lemma}[Adaptation by a retraction]
Suppose $q:G\to Q$ admits a group-homomorphic section $s$ with image $H$,
and $p:G\to L$ is a homomorphism. For a normalized additive three-cocycle
$v$ on $L$, the cochain
\[
\omega=p^*v-(psq)^*v
\]
is a normalized three-cocycle identically zero on $H$. Thus the ordinary
subgroup algebra of $H$ exists in $\Vec_G^{\exp(2\pi i\omega)}$.
\end{lemma}
\begin{proof}
Both terms are pullbacks along homomorphisms and hence cocycles.
For inputs in $H$, the retraction $sq$ is the identity, so the terms
cancel. The ordinary subgroup multiplication is therefore associative;
the usual averaging splitting proves separability.
\end{proof}

For $\R{414}$, set
\[
G_0=S_4\times C_2,\qquad
H_0=\langle((34),0),((12),1)\rangle.
\]
The map $(\operatorname{sgn},\mathrm{id}):G_0\to C_2^2$ restricts to an
isomorphism on $H_0$. For $\R{531}$, set
\[
G_1=S_4\times_{C_2}C_4,\qquad
H_1=\langle((1324),1)\rangle.
\]
Projection to $C_4$ restricts to an isomorphism on $H_1$. In both cases,
use projection to $S_4$ for $p$, and a fourth-root cocycle $v$ transferred
from a Sylow $D_8$ subgroup of $S_4$. The resulting relative cocycle and
the trivial subgroup cochain give the required unitary categories.

The actual cocycle tables, group tables, subgroup embeddings, simple
bimodule actions and basis identifications are supplied. Both replay
implementations check all $48^4=5\,308\,416$ cocycle identities, the
subgroup algebra, and all 1728 fusion coefficients. The origin by transfer
explains discovery; existence does not depend on an unproved claim that
the chosen transfer generators exhaust cohomology.

These use standard categorical constructions. The specific cocycles and
census identifications are recorded here without a claim of novelty up
to tensor equivalence.

\section{A general-complex exclusion for \texorpdfstring{$\R{350}$}{R12-1-350}}
\begin{theorem}
The ring $\R{350}$ admits no complex fusion categorification.
\end{theorem}
\begin{proof}[Proof with finite exact certificate]
Its basis is $g,g\rho$ for $g\in S_3$, where $\rho$ is central and
\[
\rho^2=1+\left(\sum_{g\in S_3}g\right)\rho,\qquad
d=\FPdim(\rho)=3+\sqrt{10}.
\]
Published Lemma 3.1 of Schopieray \cite{Sch} gives a pseudo-unitary Galois
conjugate of any categorification whose simple FP dimensions lie in
$\{1,d\}$. It does not assume commutativity, rank six, or an initial
pivotal structure. We may therefore work with a positive spherical
categorification, and write $D=120+36\sqrt{10}$.

The six irreducible representations of the complex based algebra have
degrees $(1,1,1,1,2,2)$ and formal codegrees
\[
D,\ D'=120-36\sqrt{10},\ 12,\ 12,\ 6,\ 6.
\]
Accordingly, the six distinct central simples in $I(1)$ have dimensions
$1,d^2,1+3d,1+3d,2+6d,2+6d$ and unit restriction coefficients
$1,1,1,1,2,2$. Their noninvertible restriction coefficients are constant
on the three conjugacy classes of sizes $2,3,1$.

The exact induction Gram matrix is
$B_{ij}=\sum_{a,k}N_{ai}^{k}N_{ka^*}^{j}$.
Its unit column has value twelve at the unit, zero at other invertibles,
and six at each noninvertible. The dimension and unit-column conditions
leave exactly thirty ordered low restriction frames. Twenty-six have
an exact negative Gram witness. Four survive.

For each survivor, possible remaining restriction columns are bounded
using nonnegativity, the full Gram kernel, the rank-one PSD condition,
the $d$-number condition on their central dimensions $z$, and algebraic integrality
of $D/z$. This yields respectively $30,43,43,28$ columns. Centrality
justifies compression to the five nonunit conjugacy classes; no
unproved invariance of arbitrary noncentral objects is assumed.

The categorical Cauchy theorem restricts twist-exponent primes to
$2,3,5$, the prime divisors of $DD'=1440$. The dimension field has
conductor forty. The first six normalized relative traces of all allowed
roots of unity form a finite set of 45 tuples. Orders dividing 7200,
together with the all-zero tuple for higher prime-power layers, cover
this entire trace image. This is not a bound imposed on the true exponent.

Duality and adjoint parity fix the second indicators of noninvertibles. The three
invertible transpositions retain all eight independent sign choices,
even if additional categorical constraints could relate them. The
Ng--Schauenburg center trace identities \cite{NS} and the rotation-space bounds
\[
-\frac{[X^j:1]}{p-1}\ \leq\ \overline{\nu_j(X)}\ \leq\ [X^j:1],
\quad j=3,4,5,6,
\]
where $p$ is the least prime divisor of $j$, give an exact rational
nonnegative linear system for every surviving frame and sign assignment.

For each of these $4\cdot8=32$ systems $Ax=b$, the archive supplies an
integer vector $y$ with $A^{\mathsf T}y\geq0$ and $b^{\mathsf T}y<0$.
Thus no nonnegative solution exists. An independent implementation
reconstructs every system and checks every inequality using integers and
rational arithmetic, with no optimizer or imported rank-six code.
The initial Galois reduction proves the general-complex conclusion.
\end{proof}

\section{What remains open, and what was learned}
\subsection{Dimension 32: eighteen rings}
The Ising--pointed construction resolves two cases. Checks of actual
soft tensor actions on the Choi--Lu--Sun dimension-16 categories recover
known controls, not further targets. Their coherence was checked, but
the restricted tensorator ansatz is not a classification of all actions.
A possible common isotropic involution reduction does not yet exhaust
all Morita sources: the reverse involution can fix noninvertibles and
can lie inside their stabilizer. These branches remain in the problem.

\subsection{Dimension 48: twenty rings}
The untwisted group-theoretical search was expanded to all 52 groups of
order 48 and all 1722 conjugacy classes of subgroups. Subgroup-cochain
twists and additional relative cocycles were examined. The two models
above are the accepted new outputs. This is not a proof that all possible
associators were exhausted, nor a theorem that every hypothetical
dimension-48 categorification is group-theoretical. The outstanding
cohomology and categorical converse gaps are recorded explicitly.

\subsection{Dimension 144: one Morita branch remains}
All four rings $\R{180},\R{181},\R{320},\R{321}$ remain unresolved.
The first actual $1+z$ condensation has integral type
\[
1^6,\ 3^2,\ 2^{4t},\ 4^{3-t},\qquad 0\leq t\leq3.
\]
Its pointed group is $S_3$. The normal $C_3$ cannot fix a dimension-two
or dimension-four simple, since the corresponding internal end would
have nonintegral dimension $4/3$ or $16/3$. Orbit divisibility forces
$t=0$ or $3$.

The $t=0$ branch is now excluded. A second actual $C_3$ condensation
would give integral type $(1^{24},4^3)$ and four odd simples of dimension
$3\sqrt2$. A formal-codegree argument forces a dimension-24 near-group
subcategory of type $\mathrm{NG}(H_8,2)$. The rank-four odd module has
normalized dimensions $\sqrt6$. The applicable near-group classification
is legitimate here: dimension 24 forces solvability and hence
unitarizability \cite{ENO,GHR,IT}.

The group-theoretical module formula forces the untwisted $D_8$
near-group case, with restricted module dual $\Rep(S_4)$. Both subgroup
Schur-cochain classes are retained. Clifford theory \cite{Gal} then
makes its full module dual a $C_3$-equivariantization of $\Rep(S_4)$.
Every such action fixes the five simple isomorphism classes, giving type
$(1^6,2^3,3^6)$. This contradicts the required dual type $(1^{24},4^3)$.

Consequently the unique possible first-dual type is
\[
1^6,\ 2^{12},\ 3^2,\ (3\sqrt2)^4.
\]
The original rank-three subcategory is also proved tensor-equivalent to
ordinary $\Rep(S_3)$, so its dimension-six regular algebra exists.
A central embedding, the remaining projective stabilizer classes, and
the required categorical symmetry lifts are not yet established.
Eliminating one Morita branch is not an exclusion of the four rings.

\subsection{The four remaining quadratic rings}
The unresolved quadratic IDs are $\R{332}$, $\R{561}$, $\R{623}$ and
$\R{628}$. For the two commutative cases $\R{561}$ and $\R{628}$, a
central Tannakian $C_3$ lift would yield forbidden rank-four quotients.
This rules out that lift; it does not rule out the original rings.
For $\R{623}$, the higher-indicator calculation forces second indicator
$-1$ for its pointed involution in any pseudo-unitary realization, so the
ordinary algebra $1+z$ cannot supply a model there. The opposite-sign
control agrees with the known $\R{326}$ category. Surviving indicator
branches for $\R{623}$ remain.
Neither feasible moment data nor failure to find a coherent action has
been treated as a categorical model or exclusion.

\section{Evidence, replay, and publication boundary}
The archive contains the full 632-record ledger, hashes, new decisions,
the 46 unresolved tensors, detailed family proofs, exact model data,
integer exclusion rays, independent replay programs and their receipts.
Every accepted new model is checked against all 1728 input coefficients.
The principal replay programs are tested in normal and optimized Python;
their checks do not disappear when assertions are disabled.

The finite replays verify the encoded arithmetic. The categorical
converses, scope reductions and coverage arguments are supplied as
mathematical proofs and were separately reviewed within the project.
That distinction is essential: a correct scalar computation cannot
replace a missing categorical premise.

The GitHub repositories and public website have not been revised. The
prior round remains a frozen baseline; this package is a new offline
continuation prepared for review.

\clearpage
\begin{thebibliography}{9}
\bibitem{EGNO} P. Etingof, S. Gelaki, D. Nikshych and V. Ostrik,
\emph{Tensor Categories}, AMS, 2015. Sections 7.4 and 9.7.
\href{https://math.mit.edu/~etingof/egnobookfinal.pdf}{Author manuscript}.
\bibitem{Sch} A. Schopieray, \emph{Categorification of integral group rings
extended by one dimension}, J. London Math. Soc. 108 (2023), 1617--1641.
\href{https://doi.org/10.1112/jlms.12789}{Published Lemma 3.1}.
\bibitem{ENO} P. Etingof, D. Nikshych and V. Ostrik,
\emph{Weakly group-theoretical and solvable fusion categories},
Adv. Math. 226 (2011), 176--205.
\href{https://arxiv.org/abs/0809.3031}{arXiv:0809.3031}, Theorem 1.6.
\bibitem{GHR} C. Galindo, S.-M. Hong and E. C. Rowell,
\emph{Generalized and quasi-localizations of braid group representations},
IMRN 2013, 693--731.
\href{https://arxiv.org/abs/1105.5048}{arXiv:1105.5048}, Theorem 5.20.
\bibitem{IT} M. Izumi and H. Tucker,
\emph{Algebraic realization of noncommutative near-group fusion categories},
\href{https://arxiv.org/abs/1908.01655}{arXiv:1908.01655},
Theorems 1.2 and 3.2.
\bibitem{Gal} C. Galindo, \emph{Clifford theory for graded fusion categories},
\href{https://arxiv.org/abs/1010.5283}{arXiv:1010.5283v3}, Corollary 6.4.
\bibitem{Ost} V. Ostrik, \emph{On formal codegrees of fusion categories},
\href{https://arxiv.org/abs/0810.3242}{arXiv:0810.3242}, Theorems 1.2 and 1.8.
\bibitem{BNRW} P. Bruillard, S.-H. Ng, E. C. Rowell and Z. Wang,
\emph{Rank-finiteness for modular categories},
\href{https://arxiv.org/abs/1310.7050}{arXiv:1310.7050}, Theorem 3.9.
\bibitem{NS} S.-H. Ng and P. Schauenburg,
\emph{Frobenius--Schur indicators and exponents of spherical categories},
\href{https://arxiv.org/abs/math/0601012}{arXiv:math/0601012}.
\end{thebibliography}
\end{document}
